\documentclass{article}
\usepackage[T1]{fontenc}
\usepackage{lmodern}
\usepackage{graphicx}
\usepackage{amsmath,amssymb}
\usepackage[a4paper,margin=2cm]{geometry}
\usepackage[absolute,overlay]{textpos}
\setlength{\TPHorizModule}{1mm}
\setlength{\TPVertModule}{1mm}
\usepackage[indonesian]{babel}
\usepackage{hyperref}
\setlength{\emergencystretch}{2em}
% unit: Halaman 43

\begin{document}

% === Halaman 43 (python/native) ===

43

• RC4 membangkitkan kunci-alir (keystream) dalam satuan byte setiap 

kalinya, yang kemudian di-XOR-kan dengan byte plainteks (operasi 

bitwise XOR)

• Untuk membangkitkan kunci-alir, cipher menggunakan status internal 

yang terdiri dari:

• Larik S0, S1, …, S255. Elemen-elemen di dalam larik S dipermutasi berdasarkan 

kunci rahasia K dengan panjang variabel. 

• Dua buah pencacah indeks, i dan j  

• RC4 terdiri dari dua sub-proses:


1. Key-Scheduling Algorithm (KSA)


2. Pseudo-random generation algorithm (PRGA).

\end{document}
